\documentclass[11pt]{article}
\usepackage[margin=1in]{geometry}
\usepackage{amsmath, amssymb, bm}
\usepackage{algorithm, algpseudocode}
\usepackage{newtxtext,newtxmath}
\usepackage{hyperref}

\title{Adjoint Sensitivity Formulation for WaterLily.jl}
\author{}
\date{}

\newcommand{\vect}[1]{\bm{#1}}
\newcommand{\dd}{\,\mathrm{d}}
\newcommand{\pp}{\partial}

\begin{document}
\maketitle

%----------------------------------------------------------------------
\section{Governing Equations}
%----------------------------------------------------------------------

WaterLily.jl solves the incompressible Navier--Stokes equations on a
Cartesian grid with immersed boundaries modelled via the Boundary Data
Immersion Method (BDIM)~\cite{Maertens2015}.

\subsection{Continuous equations}

The incompressible Navier--Stokes equations in non-dimensional form are
\begin{align}
  \frac{\pp \vect{u}}{\pp t} + (\vect{u}\cdot\nabla)\vect{u}
    &= -\nabla p + \nu\,\nabla^{2}\vect{u}, \label{eq:momentum}\\
  \nabla\cdot\vect{u} &= 0. \label{eq:continuity}
\end{align}
Here $\vect{u}$ is the velocity vector, $p$ is the kinematic pressure, and
$\nu$ is the kinematic viscosity.

\subsection{BDIM immersed boundary treatment}

Let $d(\vect{x})$ denote the signed distance from $\vect{x}$ to the body
surface ($d>0$ in fluid, $d<0$ inside body), and define the zeroth-moment
kernel
\begin{equation}\label{eq:mu0}
  \mu_0(d;\epsilon) = \operatorname{kern}_0\!\bigl(\operatorname{clamp}(d/\epsilon, -1, 1)\bigr),
  \qquad
  \operatorname{kern}_0(\hat d) = \tfrac{1}{2} + \tfrac{1}{2}\hat d
    + \frac{\sin(\pi\hat d)}{2\pi},
\end{equation}
where $\epsilon$ is the immersion width (typically one cell).
For a static body with zero velocity, the simplified (zeroth-order) BDIM
predictor is
\begin{equation}\label{eq:bdim_pred}
  \vect{u}^{*} = \mu_0\bigl(\vect{u}^{n} + \Delta t\,\vect{f}^{n}\bigr),
\end{equation}
where $\vect{f}^{n}$ is the explicit convective--diffusive right-hand side
evaluated at the old velocity $\vect{u}^{n}$.

\subsection{Pressure projection}

After the predictor, the velocity is projected onto the divergence-free
subspace by solving a variable-coefficient Poisson equation
\begin{equation}\label{eq:poisson}
  \nabla\cdot\bigl(\mu_0\,\nabla\phi\bigr) = \nabla\cdot\vect{u}^{*}
\end{equation}
and correcting
\begin{equation}\label{eq:proj}
  \vect{u}^{n+1} = \vect{u}^{*} - \mu_0\,\nabla\phi,
  \qquad
  p^{n+1} = \phi/\Delta t.
\end{equation}
The $\mu_0$ weighting ensures that the pressure gradient vanishes inside
the body, enforcing a no-slip condition without explicit forcing.

%----------------------------------------------------------------------
\section{Discrete Operators on the Staggered Grid}
%----------------------------------------------------------------------

WaterLily uses a staggered (MAC) grid with ghost cells.  For a domain of
$N_1\times N_2$ interior cells, the arrays have size
$\bar N_k = N_k + 2$ (one ghost cell on each side).

\begin{itemize}
  \item \textbf{Pressure} $p_I$, scalar at cell centres, $I\in\bar N_1\times\bar N_2$.
  \item \textbf{Velocity} $u_{I,i}$, $i$-component at face $I$ in direction $i$.
  \item \textbf{Interior pressure cells:} $I\in[2,\bar N_1{-}1]\times[2,\bar N_2{-}1]$.
  \item \textbf{Interior velocity faces (component $i$):}
    $I_k\in\begin{cases}[3,\bar N_k{-}1]&k=i,\\[2pt][2,\bar N_k{-}1]&k\neq i.\end{cases}$
\end{itemize}

The discrete gradient operator $\partial_i$ and divergence operator $\operatorname{div}$ are
\begin{equation}
  \partial_i p_I = p_I - p_{I-\delta_i},
  \qquad
  \operatorname{div}(I, \vect{u}) = \sum_{i=1}^{D}(u_{I+\delta_i,\,i} - u_{I,i}).
\end{equation}

The convective--diffusive flux is computed using a central-difference scheme
(CDS) for the convective term and standard second-order finite differences
for diffusion:
\begin{equation}
  f_{I,i} = -\sum_j\bigl[\phi^u_{j}(I,i) - \nu\,\partial_j u_{I,i}\bigr]
    \quad\text{(conservative flux form)},
\end{equation}
where $\phi^u_j$ is the CDS flux $\phi^u = \bar u\cdot\tfrac{1}{2}(u_L+u_R)$.

%----------------------------------------------------------------------
\section{Boundary Conditions}
%----------------------------------------------------------------------

\begin{itemize}
  \item \textbf{Velocity, normal component (Dirichlet):}
    $u_{I,i} = U_i^{\mathrm{BC}}$ at boundary faces in direction $i$.
  \item \textbf{Velocity, tangential component (Neumann):}
    $u_{\text{ghost},i} = u_{\text{interior},i}$ (zero-gradient).
  \item \textbf{Pressure (Neumann):}
    Implicitly enforced by setting $\mu_0=0$ at domain boundaries,
    so $\mu_0\,\partial_i p = 0$ there.
  \item \textbf{Ghost pressure cells:}
    Not modified by the solver; values persist from the previous step.
\end{itemize}

%----------------------------------------------------------------------
\section{Primal (Forward) Solver Pseudo-Code}
\label{sec:forward}

The \emph{forward} (or \emph{primal}) solver computes the physical flow
field by marching the Navier--Stokes equations forward in time.  This
terminology distinguishes it from the \emph{adjoint} solver
(Section~\ref{sec:adjoint}), which propagates sensitivity information
backward through the time steps.  The analogy mirrors forward- versus
reverse-mode algorithmic differentiation~\cite{Martins2021}.
%----------------------------------------------------------------------

\begin{algorithm}[H]
\caption{One-step forward solve: \texttt{ns\_onestep!}}
\label{alg:onestep}
\begin{algorithmic}[1]
\Require Previous state $\vect{y}^n = [\vect{u}^n;\,p^n]$, design variable $c_x$
\Ensure  New state $\vect{y}^{n+1} = [\vect{u}^{n+1};\,p^{n+1}]$
\State Unpack $\vect{u}^n, p^n$ from $\vect{y}^n$
\State Construct body SDF: $d(\vect{x}) = \|\vect{x}-\vect{x}_c(c_x)\| - R$
\State Compute $\mu_0$ via \texttt{measure!} using Eq.~\eqref{eq:mu0}
\State Build multigrid Poisson solver with coefficient $L=\mu_0$
\Statex
\Comment{\textit{BDIM predictor}}
\State Compute explicit RHS: $\vect{f}^n \gets \texttt{conv\_diff!}(\vect{u}^n)$ \Comment{CDS convection + diffusion}
\ForAll{interior faces $(I,i)$}
  \State $u^*_{I,i} \gets \mu_{0,I,i}\,(u^n_{I,i} + \Delta t\,f^n_{I,i})$
  \Comment{Eq.~\eqref{eq:bdim_pred}}
\EndFor
\State Apply velocity BCs: \texttt{BC!}$(\vect{u}^*,\,\vect{U}^{\mathrm{BC}})$
\Statex
\Comment{\textit{Pressure projection}}
\ForAll{interior pressure cells $I$}
  \State $\sigma_I \gets \operatorname{div}(I, \vect{u}^*)$
  \Comment{Divergence source}
\EndFor
\State $\phi \gets \texttt{solver!}(\nabla\cdot(\mu_0\nabla\phi)=\sigma;\;\mathrm{tol}=\varepsilon_{\mathrm{mach}})$
  \Comment{Tight multigrid V-cycle}
\ForAll{interior pressure cells $I$, directions $i$}
  \State $u^{n+1}_{I,i} \gets u^*_{I,i} - \mu_{0,I,i}\,\partial_i\phi_I$
  \Comment{Eq.~\eqref{eq:proj}}
\EndFor
\State $p \gets \phi / \Delta t$
\State Apply velocity BCs: \texttt{BC!}$(\vect{u}^{n+1},\,\vect{U}^{\mathrm{BC}})$
\Statex
\Comment{\textit{Pressure pinning (removes constant null space)}}
\State $p_{\mathrm{ref}} \gets p_{I_{\mathrm{pin}}}$
\ForAll{interior pressure cells $I$}
  \State $p_I \gets p_I - p_{\mathrm{ref}}$
\EndFor
\State Pack $\vect{y}^{n+1} \gets [\vect{u}^{n+1};\,p^{n+1}]$
\end{algorithmic}
\end{algorithm}

%----------------------------------------------------------------------
\section{Residual Formulation for the Adjoint}
%----------------------------------------------------------------------

The adjoint method requires a residual function
$\vect{R}(\vect{y}^{n+1}, \vect{y}^{n}, \vect{x}_d)=\vect{0}$
that is satisfied at the forward solution, following the discrete approach
of Martins \& Ning~\cite{Martins2021}, Eq.~(6.37).

\subsection{State vector layout}

The state vector $\vect{y}\in\mathbb{R}^{n_y}$ concatenates the full
velocity and pressure arrays, \emph{including ghost cells}:
\begin{equation}\label{eq:state_vec}
  \vect{y}
  = \begin{bmatrix} \operatorname{vec}(\vect{u}) \\ \operatorname{vec}(p) \end{bmatrix}
  = \begin{bmatrix}
      \operatorname{vec}(u_{\cdot,1}) \\[2pt]
      \operatorname{vec}(u_{\cdot,2}) \\[2pt]
      \vdots \\[2pt]
      \operatorname{vec}(u_{\cdot,D}) \\[2pt]
      \operatorname{vec}(p)
    \end{bmatrix}
  \;\in\mathbb{R}^{n_y},
  \qquad
  n_y = (D+1)\,N_p,
  \qquad
  N_p = \prod_{k=1}^{D}\bar N_k.
\end{equation}
Here $N_p=\bar N_1 \bar N_2$ is the number of scalar DOFs per field
(including ghosts), $N_u = D\,N_p$ is the total velocity DOF count,
and $\operatorname{vec}(\cdot)$ denotes column-major flattening.

\paragraph{Memory layout.}
WaterLily stores $\vect{u}$ as a $(D{+}1)$-dimensional array of size
$(\bar N_1, \bar N_2, D)$ and $p$ as a $D$-dimensional array of size
$(\bar N_1, \bar N_2)$.  In Julia's column-major order the velocity
component $u_{I,i}$ at multi-index $I=(I_1,I_2)$ maps to the linear
index
\begin{equation}\label{eq:lin_vel}
  \ell^{u}(I, i) = (i-1)\,N_p + I_1 + (I_2-1)\,\bar N_1,
  \qquad i=1,\ldots,D,
\end{equation}
and the pressure $p_I$ maps to
\begin{equation}\label{eq:lin_prs}
  \ell^{p}(I) = N_u + I_1 + (I_2-1)\,\bar N_1.
\end{equation}
Equivalently, $y[\ell^{u}(I,i)] = u_{I,i}$ and $y[\ell^{p}(I)] = p_I$.

\paragraph{DOF classification.}
Each entry of $\vect{y}$ falls into one of four categories, which
determines which residual equation governs it:

\begin{center}
\renewcommand{\arraystretch}{1.3}
\begin{tabular}{llll}
  \hline
  \textbf{Block} & \textbf{Index range} & \textbf{Interior set}
    & \textbf{Residual} \\
  \hline
  Velocity $u_{I,i}$
    & $\ell^{u}(I,i)\in[1,N_u]$
    & $I\in\texttt{inside\_u}(\bar N,i)$
    & Eq.~\eqref{eq:res_mom} \\
  \quad (ghost/BC)
    &
    & complement
    & Eq.~\eqref{eq:res_bc_vel} \\
  Pressure $p_I$
    & $\ell^{p}(I)\in[N_u{+}1,\,n_y]$
    & $I\in\texttt{inside}(p)$
    & Eq.~\eqref{eq:res_cont} or \eqref{eq:res_pin} \\
  \quad (ghost)
    &
    & complement
    & Eq.~\eqref{eq:res_bc_prs} \\
  \hline
\end{tabular}
\end{center}

\paragraph{Concrete example ($D=2$, $N=(32,16)$).}
With $\bar N = (34, 18)$: $N_p = 612$, $N_u = 1224$, $n_y = 1836$.
The state vector is partitioned as
\[
  \underbrace{y_1,\ldots,y_{612}}_{u_{x}}
  \;\;
  \underbrace{y_{613},\ldots,y_{1224}}_{u_{y}}
  \;\;
  \underbrace{y_{1225},\ldots,y_{1836}}_{p}.
\]

\subsection{Residual equations}

The residual vector $\vect{R}\in\mathbb{R}^{n_y}$ is indexed identically
to $\vect{y}$: entry $R[\ell^{u}(I,i)]$ is the equation that governs
$u_{I,i}$, and $R[\ell^{p}(I)]$ governs $p_I$.  The residual decomposes
into five blocks, each zero at the forward solution.

\subsubsection{Interior momentum}

For each component $i$ and every interior face
$I\in\texttt{inside\_u}(\bar N, i)$:
\begin{equation}\label{eq:res_mom}
  \boxed{%
  R[\ell^{u}(I,i)]
  = u^{n+1}_{I,i}
  - \underbrace{\mu_0\bigl(u^n_{I,i} + \Delta t\,f^n_{I,i}\bigr)}_{\text{BDIM predictor}}
  + \underbrace{\Delta t\,\mu_0\,\partial_i p^{n+1}_I}_{\text{pressure correction}}.}
\end{equation}
Here $\mu_0 = \mu_0\bigl(d(\vect{x}_{\mathrm{face}};\,c_x),\,\epsilon\bigr)$
is evaluated at the face location and depends on $c_x$ through the signed
distance~$d$.  The explicit flux $f^n_{I,i}$ depends only on
$\vect{u}^{n}$ (the previous step), so the residual is linear in
$\vect{u}^{n+1}$ and $p^{n+1}$.

Expanding the pressure gradient:
\begin{equation}
  \Delta t\,\mu_0\,\partial_i p^{n+1}_I
  = \Delta t\,\mu_0\bigl(p^{n+1}_I - p^{n+1}_{I-\delta_i}\bigr).
\end{equation}

\subsubsection{Interior continuity}

For every interior pressure cell $I\in\texttt{inside}(p)$ that is
\emph{not} pinned (see below):
\begin{equation}\label{eq:res_cont}
  \boxed{%
  R[\ell^{p}(I)] = \operatorname{div}(I,\,\vect{u}^{n+1})
  = \sum_{i=1}^{D}\bigl(u^{n+1}_{I+\delta_i,\,i} - u^{n+1}_{I,i}\bigr).}
\end{equation}
This equation involves only velocity, not pressure.

\subsubsection{Pressure pinning}

Two classes of interior pressure cells are \emph{pinned}
($R[\ell^{p}(I)] = p^{n+1}_I$) instead of satisfying
the continuity equation:
\begin{equation}\label{eq:res_pin}
  \boxed{R[\ell^{p}(I)] = p^{n+1}_{I}}
  \qquad\text{if}\quad
  I = I_{\mathrm{pin}}
  \;\;\text{or}\;\;
  d\bigl(\vect{x}_{\mathrm{center}}(I)\bigr) < -\bigl(\epsilon + \tfrac{1}{2}\bigr).
\end{equation}

\begin{itemize}
  \item \textbf{Reference pin} ($I_{\mathrm{pin}}$, one cell):
    Removes the constant null space of the discrete Laplacian under
    all-Neumann pressure boundary conditions.  Since only pressure
    \emph{gradients} appear in the momentum equations, a global
    pressure shift does not affect the velocity field.  Pinning one
    cell to zero selects the unique pressure that satisfies the
    residual.

  \item \textbf{Body-interior pins} (cells deep inside the body):
    When $d < -(\epsilon+\frac{1}{2})$ at the cell centre, every
    surrounding face has $d_{\mathrm{face}} < -\epsilon$, hence
    $\mu_0 = \operatorname{kern}_0(-1) = 0$.  Pressure at these cells
    does not appear in \emph{any} momentum equation (the gradient
    coefficient $\Delta t\,\mu_0 = 0$ on all adjacent faces).
    Without pinning, the corresponding columns of
    $\pp\vect{R}/\pp\vect{y}$ are structurally zero and the Jacobian
    is singular.
\end{itemize}

\subsubsection{Ghost/boundary velocity}

All velocity DOFs \emph{outside} the interior set satisfy a boundary
condition constraint:
\begin{equation}\label{eq:res_bc_vel}
  \boxed{R[\ell^{u}(I,i)] = u^{n+1}_{I,i} - \texttt{BC}\bigl(\vect{u}^{n+1}\bigr)_{I,i},}
\end{equation}
where $\texttt{BC}(\cdot)$ applies WaterLily's boundary condition
operator to a copy of $\vect{u}^{n+1}$ and returns the expected
ghost-cell value.  This captures both types of boundary condition:
\begin{itemize}
  \item \textbf{Dirichlet} (normal component, $i=j$ at boundary in
    direction~$j$): $\texttt{BC}(\vect{u})_{I,i} = U_i^{\mathrm{BC}}$
    (constant), so
    $R = u^{n+1}_{I,i} - U_i^{\mathrm{BC}}$.
  \item \textbf{Neumann} (tangential component, $i \neq j$):
    $\texttt{BC}(\vect{u})_{I,i} = u^{n+1}_{I\pm\delta_j,\,i}$
    (the interior neighbour), so
    $R = u^{n+1}_{I,i} - u^{n+1}_{I\pm\delta_j,\,i}$.
\end{itemize}

\subsubsection{Ghost pressure}

Pressure DOFs outside the interior set satisfy an identity constraint:
\begin{equation}\label{eq:res_bc_prs}
  \boxed{R[\ell^{p}(I)] = p^{n+1}_I - p^{n}_I.}
\end{equation}
Ghost pressure cells are not modified by the Poisson solver (boundary
$\mu_0=0$ decouples them from the interior), and the pressure pinning
shift is applied only to interior cells, so $p^{n+1}_I = p^{n}_I$ and
$R=0$ at the forward solution.

\subsection{Residual vector summary}

Collecting the five blocks, the full residual vector reads
\begin{equation}\label{eq:res_full}
  \vect{R}(\vect{y}^{n+1},\vect{y}^{n},\vect{x}_d) =
  \begin{bmatrix}
    R[\ell^{u}(I,i)] \\[6pt] R[\ell^{p}(I)]
  \end{bmatrix}
  =
  \begin{cases}
    u^{n+1}_{I,i} - \mu_0(u^n_{I,i}+\Delta t\,f^n_{I,i})
      + \Delta t\,\mu_0\,\partial_i p^{n+1}_I
      & I\in\mathcal{U}_{\mathrm{int}}(i), \\[4pt]
    u^{n+1}_{I,i} - \texttt{BC}(\vect{u}^{n+1})_{I,i}
      & I\notin\mathcal{U}_{\mathrm{int}}(i), \\[8pt]
    \operatorname{div}(I,\vect{u}^{n+1})
      & I\in\mathcal{P}_{\mathrm{int}}\setminus\mathcal{P}_{\mathrm{pin}}, \\[4pt]
    p^{n+1}_I
      & I\in\mathcal{P}_{\mathrm{pin}}, \\[4pt]
    p^{n+1}_I - p^{n}_I
      & I\notin\mathcal{P}_{\mathrm{int}},
  \end{cases}
\end{equation}
where
$\mathcal{U}_{\mathrm{int}}(i) = \texttt{inside\_u}(\bar N,i)$,\;
$\mathcal{P}_{\mathrm{int}} = \texttt{inside}(p)$,\;
and
$\mathcal{P}_{\mathrm{pin}} = \{I_{\mathrm{pin}}\}
  \cup \{I : d(\vect{x}_{\mathrm{center}}(I)) < -(\epsilon+\tfrac{1}{2})\}$.

%----------------------------------------------------------------------
\section{Jacobian Structure}
\label{sec:jacobian}
%----------------------------------------------------------------------

The Jacobian $\vect{J} = \pp\vect{R}/\pp\vect{y}^{n+1}$ has a
$2\times2$ block structure when partitioned into velocity and pressure DOFs:
\begin{equation}\label{eq:jacobian}
  \vect{J} =
  \begin{bmatrix}
    \vect{J}_{uu} & \vect{J}_{up} \\[4pt]
    \vect{J}_{pu} & \vect{J}_{pp}
  \end{bmatrix}.
\end{equation}
The entries of each block follow directly from the
residual~\eqref{eq:res_full}:

\paragraph{$\vect{J}_{uu}$ (velocity--velocity).}
Interior momentum rows have
$\pp R^{\mathrm{mom}}_{I,i}/\pp u^{n+1}_{I,i} = 1$ on the diagonal;
all off-diagonal entries are zero because the predictor depends on
$\vect{u}^n$, not $\vect{u}^{n+1}$.
Ghost Dirichlet rows contribute $+1$ on the diagonal.
Ghost Neumann rows have $+1$ on the diagonal and $-1$ at the interior
neighbour.  Hence $\vect{J}_{uu}\approx\vect{I}$.

\paragraph{$\vect{J}_{up}$ (velocity--pressure).}
Interior momentum rows have two entries per row from the pressure gradient:
\[
  \frac{\pp R^{\mathrm{mom}}_{I,i}}{\pp p^{n+1}_I} = +\Delta t\,\mu_0,
  \qquad
  \frac{\pp R^{\mathrm{mom}}_{I,i}}{\pp p^{n+1}_{I-\delta_i}} = -\Delta t\,\mu_0.
\]
Ghost velocity rows contribute zero (boundary conditions do not depend on
pressure).  Thus $\vect{J}_{up} = \Delta t\,\mu_0\,\vect{G}$ where
$\vect{G}$ is the discrete gradient operator.

\paragraph{$\vect{J}_{pu}$ (pressure--velocity).}
Continuity rows contribute the discrete divergence stencil ($\pm 1$
entries).  Pinned pressure rows and ghost pressure rows contribute zero.
Thus $\vect{J}_{pu} = \vect{D}$ (with rows zeroed for pinned cells),
where $\vect{D}$ is the discrete divergence operator.

\paragraph{$\vect{J}_{pp}$ (pressure--pressure).}
Continuity rows have no pressure dependence: $\vect{0}$.
Pinned rows (reference pin + body-interior) have $+1$ on the diagonal.
Ghost pressure identity rows have $+1$ on the diagonal.

\paragraph{Full structure.}
Ordering DOFs as
(ghost $u$, interior $u$, ghost $p$, pinned $p$, interior $p$):
\begin{equation}
  \vect{J} \sim
  \begin{bmatrix}
    \vect{I}_{\mathrm{bc}} & * & \vect{0} & \vect{0} & \vect{0} \\
    \vect{0} & \vect{I} & \vect{0} & \Delta t\,\mu_0\,\vect{G}_{\mathrm{pin}} & \Delta t\,\mu_0\,\vect{G}_{\mathrm{int}} \\
    \vect{0} & \vect{0} & \vect{I} & \vect{0} & \vect{0} \\
    \vect{0} & \vect{0} & \vect{0} & \vect{I} & \vect{0} \\
    \vect{0} & \vect{D} & \vect{0} & \vect{0} & \vect{0}
  \end{bmatrix},
\end{equation}
where $\vect{I}_{\mathrm{bc}}$ has $+1$ on the diagonal with possible
$-1$ off-diagonal entries from Neumann coupling, and $*$ denotes the
Neumann coupling between ghost and interior velocity.

The Schur complement for interior pressure, after eliminating velocity, is
\begin{equation}
  \vect{S} = -\vect{D}\,\vect{J}_{uu}^{-1}\,\Delta t\,\mu_0\,\vect{G}
  \approx -\Delta t\,\nabla\cdot(\mu_0\,\nabla),
\end{equation}
which is the (negative) discrete weighted Laplacian.  Without pinning,
$\vect{S}$ has a one-dimensional null space (constant pressure mode) plus
additional null dimensions from body-interior cells.  The pinning
equations replace the corresponding rows with $\vect{e}_I^{\top}$,
lifting the rank deficiency.

%----------------------------------------------------------------------
\section{Adjoint Derivative Computation}
\label{sec:adjoint}
%----------------------------------------------------------------------

Following Martins \& Ning~\cite{Martins2021}, \S6.7, the total derivative
of a scalar objective $f(\vect{y},\vect{x}_d)$ with respect to the design
variables is
\begin{equation}\label{eq:total_deriv}
  \frac{\dd f}{\dd \vect{x}_d}
  = \frac{\pp f}{\pp \vect{x}_d}
  - \bm{\psi}^{\!\top}\frac{\pp\vect{R}}{\pp\vect{x}_d},
\end{equation}
where the adjoint vector $\bm\psi$ satisfies the adjoint equation
\begin{equation}\label{eq:adjoint_eq}
  \vect{J}^{\top}\bm{\psi}
  = \left(\frac{\pp f}{\pp\vect{y}}\right)^{\!\top}.
\end{equation}

\subsection{Unsteady adjoint via ImplicitAD.jl}

For unsteady problems with $N_t$ time steps, ImplicitAD's
\texttt{implicit\_unsteady} applies the discrete adjoint method
time-step by time-step.  At each step $n$:

\begin{enumerate}
  \item Compute $\pp\vect{R}/\pp\vect{y}^{n+1}$ and
    $\pp\vect{R}/\pp\vect{x}_d$ via forward-mode AD (ForwardDiff.jl).
  \item Solve the adjoint linear system~\eqref{eq:adjoint_eq}.
  \item Accumulate the sensitivity contribution~\eqref{eq:total_deriv}.
\end{enumerate}

The cost is independent of the number of design variables $n_x$ (one adjoint
solve per function of interest), making it efficient for shape optimisation
where $n_x\gg n_f$~\cite{Martins2021}.

%----------------------------------------------------------------------
\section{Verification Strategy}
%----------------------------------------------------------------------

\begin{enumerate}
  \item \textbf{Residual check:}
    Evaluate $\vect{R}(\vect{y}_{\mathrm{fwd}})$ at the forward solution
    and verify $\|\vect{R}\|<\varepsilon$ (machine-precision for momentum
    and ghost cells; Poisson-tolerance-limited for continuity).
  \item \textbf{Adjoint gradient:}
    Compute $\dd f/\dd c_x$ via ImplicitAD + ForwardDiff.
  \item \textbf{Finite-difference verification:}
    Central differences
    $\dd f/\dd c_x \approx [f(c_x{+}h)-f(c_x{-}h)]/(2h)$
    and check relative agreement $<5\%$.
\end{enumerate}

%----------------------------------------------------------------------
\begin{thebibliography}{9}
\bibitem{Maertens2015}
  A.~P.~Maertens and G.~D.~Weymouth,
  ``Accurate Cartesian-grid simulations of near-body flows at intermediate
  Reynolds numbers,''
  \emph{Comput.\ Methods Appl.\ Mech.\ Engrg.}, vol.~283, pp.~106--129, 2015.
  \texttt{doi:10.1016/j.cma.2014.09.007}.

\bibitem{Martins2021}
  J.~R.~R.~A.~Martins and A.~Ning,
  \emph{Engineering Design Optimization}.
  Cambridge University Press, 2021, Ch.~6.
\end{thebibliography}

\end{document}
